% LỜI TỰA --- Quyển V
\chapter*{Lời Tựa}
\addcontentsline{toc}{chapter}{Lời Tựa}
\markboth{Lời Tựa}{Lời Tựa}

\begin{truyen}{Lời Tựa}{Two Sides of Every Coin}
\chuhoa{G}{iáo viên bắt đầu bộ sách này với một câu hỏi đơn giản: làm thế nào để mỗi buổi học có thể để lại thứ gì đó đáng nhớ hơn là bài thi?}
\textit{(A teacher began this series with a simple question: how can each class leave behind something more memorable than a test?)}

Không có điều gì là hoàn toàn tốt hay hoàn toàn xấu~--- điều quan trọng là ta nhìn nó từ góc nào.
\textit{(Nothing is entirely good or entirely bad~--- what matters is the angle from which we look at it.)}

Quyển năm đặt ra câu hỏi về hai mặt của mọi thứ. Mỗi câu chuyện được cặp đôi với một câu chuyện đối lập~--- để người đọc thấy rằng sự thật thường nằm ở giữa, không phải ở một phía. Tư duy biện chứng không phải là kỹ năng của người lớn~--- ngay từ trong trường học, học sinh đã cần học nhìn vấn đề từ nhiều chiều.
\textit{(Volume five raises the question of the two sides of everything. Each story is paired with an opposing story~--- so readers see that the truth often lies in the middle, not on one side. Dialectical thinking is not just an adult skill~--- from school onwards, students need to learn to see problems from multiple angles.)}

Bộ sách <<Truyện Tích Cảm Hứng Song Ngữ>> gồm 28 quyển, mỗi quyển một chủ đề riêng. Quyển V này mang chủ đề: \textbf{tư duy biện chứng, âm dương, mâu thuẫn và cân bằng}.
\textit{(The series ``Inspiring Bilingual Tales'' contains 28 volumes, each with its own theme. Volume V focuses on: \textbf{tư duy biện chứng, âm dương, mâu thuẫn và cân bằng}.})
\end{truyen}

\begin{baihoc}
Người khôn ngoan không chọn một mặt~--- họ hiểu cả hai và chọn điều cân bằng.
\textit{(The wise do not choose one side~--- they understand both and choose the balance.)}
\end{baihoc}

\begin{tcolorbox}[
  colback=truyenblue!6, colframe=truyenblue,
  coltitle=white, fonttitle=\bfseries,
  title={Easy English},
  arc=4pt, boxrule=1pt, breakable, enhanced
]
\textbf{Short English to remember:}

Every coin has two sides: look at both.

Balance comes from understanding opposites.
\end{tcolorbox}

\begin{tcolorbox}[
  colback=truyengold!10, colframe=truyengold!70!black,
  coltitle=black, fonttitle=\bfseries,
  title={T\textit{ự} H\textit{ọ}c Nhanh},
  arc=4pt, boxrule=1pt, breakable, enhanced
]
1. Câu chuyện nào trong quyển này khiến em thay đổi cách nhìn về một vấn đề? \textit{(Which story here changed how you see a particular issue?)}

2. Em thường nghiêng về phía nào khi xem xét một vấn đề? Lý do? \textit{(Which side do you usually lean toward when evaluating a situation? Why?)}

3. Tư duy biện chứng giúp ích gì cho em trong học tập và cuộc sống? \textit{(How does seeing both sides help you in study and in life?)}
\end{tcolorbox}

\begin{center}
\textit{\color{truyengold}``Nhìn được cả hai mặt~--- đó mới là cái nhìn trưởng thành.''}
\end{center}

\begin{flushright}
\textbf{Tôi Nguyễn Văn Sang}\\
\textit{GV THPT Nguyễn Hữu Cảnh - TP HCM}
\end{flushright}

\ngancach
